\documentclass[11pt]{article}
\usepackage[utf8]{inputenc}
\usepackage{algorithm}
\usepackage{amsmath}
\usepackage{graphicx}
\usepackage{mathtools}
\usepackage{amsfonts}
\usepackage{hyperref}
\usepackage{geometry}
\usepackage{multicol}
\usepackage{capt-of}
\usepackage{url}
\usepackage{chemmacros}
\usepackage{upgreek}
\usepackage{isomath}
\usepackage{cancel}
\usepackage{bm}
\usepackage[italian]{babel}
\geometry{a4paper, top=2.5cm, bottom=2.5cm, left=3cm, right=3cm, bindingoffset=5mm}
\usepackage[suftesi]{frontespizio}
\usepackage{hyperref}
\hypersetup{
    colorlinks=true,
    linkcolor=,
    filecolor=,      
    urlcolor=red,
    pdftitle={Scelta Bersaglio Compton},
    pdfpagemode=FullScreen,
    }
\usepackage{siunitx}
\renewcommand{\arraystretch}{1.3} % Default value: 1
\setlength{\tabcolsep}{8pt} % Default value: 6pt

\title{Scelta bersaglio}
\author{Gruppo Compton}
\date{\today}

\begin{document}

\maketitle

\begin{equation}\label{si}
    N_s = N_i \rho \frac{N_A}{M} Z dx \sigma(E)
\end{equation}
dove $\sigma$ viene valutata a $511\, \unit{KeV}$. I tre materiali presi in considerazione sono:
\begin{itemize}
    \item Alluminio;
    \item Carbonio;
    \item Rame.
\end{itemize}
A partire dai dati ottenuti attraverso il sito \href{https://physics.nist.gov/PhysRefData/Xcom/html/xcom1.html}{\underline{NIST}} abbiamo fatto le seguenti considerazioni.
Per prima cosa abbiamo escluso il rame perchè l'effetto fotoelettrico risulta di un ordine di grandezza più importante rispetto a $\mathrm{Al}$ e $\mathrm{C}$. Oltre a questo abbiamo anche calcolato la probabilità di avere effetto Compton per i tre materiali.

Utilizzando la formula (\ref{si}) abbiamo calcolato il rapporto tra il numero di fotoni deviati e il numero di fotoni incidenti. Dai risultati visibili nelle tabelle abbiamo deciso di utilizzare come scatteratore l'alluminio.

Per quanto riguarda le dimensioni del bersaglio abbiamo deciso di adottare come scatteratore un cilindro di alluminio di spessore pari a $2\, \unit{cm}$, che è circa la metà del valore di libero cammino medio dell'effetto Compton. La scelta di questa dimensione è dettata anche da due richieste: avere un rapporto $N_s/N_i$ il più alto possibile e allo stesso tempo un valore di $n_t \sigma dz << 1$ per essere in situazione di bersaglio sottile evitando così interazioni mulptiple e riducendo gli effetti di assorbimento all'interno del materiale stesso. I valori sono ripostati qui sotto:

%%tabella
%%\num{0.9e-30} \unit{\kilogram}
\begin{table}[h!]
\begin{center}
\begin{tabular}{||c c c c||} 

 \hline
 Elemento & Z & $\rho$ [\unit{\g\,\cm^{-3}}] & dx [\unit{\cm}] \\ [0.5ex] 
 \hline\hline
 Al & 13 & 2.700 & 1.914 \\ 
 \hline
 C & 6 & 2.267 & 2.005 \\
 \hline
 Cu & 29 & 8.933 & 0.60 \\
 \hline 
\end{tabular}
\caption{Bersagli disponibili in laboratorio.}
\label{table:1}

\vspace{1cm}

\begin{tabular}{||c c c c||} 

 \hline
 Elemento & $\sigma_{C}$ [\unit{\cm^{2}}] & $\sigma_{PH}$ [\unit{\cm^2}] & $\frac{\sigma_{C}}{\sigma_{TOT}}$ \\ [0.5ex] 
 \hline\hline
 Al & \num{2.860e-25} & \num{4.37e-28} & 99.19\% \\ 
 \hline
 C & \num{2.865e-25} & \num{2.15e-29} & 99.81\% \\ 
 \hline
 Cu & \num{2.848e-25} & \num{8.83e-27} & 94.67\% \\
 \hline 
\end{tabular}
\caption{Sezioni d'urto Compton (C) e Fotoelettrica (PH).}
\label{table:2}

\vspace{1cm}

\begin{tabular}{||c c c c||} 

 \hline
 Elemento & $\frac{N_s}{N_i}$ & $n_t dx \sigma$ & $\lambda_{COMP} [\unit{\cm}]$ \\ [0.5ex] 
 \hline\hline
 Al & 42.89\% & 0.4323 & 4.463 \\ 
 \hline
 C & 39.18\% & 0.392 & 5.117 \\ 
 \hline 
\end{tabular}
\caption{Considerazioni finali.}
\label{table:3}
\end{center}
\end{table}


\end{document}
